\vspace{-1mm}
\section{Preliminaries}
\vspace{-1mm}


\subsection{Problem Formulation}
\label{sec:problem_formulation}
\vspace{-1mm}

Table-text retrieval is involved from a retrieval \textit{corpus} $\mathcal{C}$, which comprises two distinct sets: a collection of passages $\mathcal{C}_P = \{P^{(1)}, \dots, P^{(n)}\}$ and a collection of tables $\mathcal{C}_T = \{T^{(1)}, \dots, T^{(m)}\}$. 
A \textit{passage} is defined as a sequence of tokens $P$, representing unstructured text. 
A \textit{table} is a structured matrix $T$, consisting of cells $T_{i, j}$, where $i$ and $j$ indicate the row index and the column index, respectively. 
Each cell $T_{i, j}$ may contain a number, date, phrase, or sentence.
We define a document as either a passage or a table.
Given a query $q$, the objective of table-text retrieval is to retrieve from corpus $\mathcal{C}$ a ranked list of documents such that the document containing the answer span $a$ is positioned among the top results.

We split a table into multiple table segments, as commonly used in existing studies. 
Since a single table can easily exceed the token limits of language models, a table $T$ is combined with its header to form a list of table segments $T = [S^{(1)}, \dots, S^{(m')}]$ (\citealt{OTT-QA}). 
This process results in (i) a corpus $\mathcal{C}$ composed of table segments $\mathcal{C}_S$ and passages $\mathcal{C}_P$ (i.e., $\mathcal{C} = \mathcal{C}_S \cup \mathcal{C}_P$) and (ii) a mapping $\mathcal{M}: \mathcal{C}_S \rightarrow \mathcal{C}_T$ to associate table segments with their original table.






\vspace{-1mm}
\subsection{Table-Text Retrieval as Graph Retrieval}
\label{sec:table_text_retrieval_as_bipartite_graph_retrieval}
\vspace{-1mm}


We adopt a graph representation, denoted by $G = (V, E, \Phi, \Gamma, \Lambda)$, to generalize various methods used in existing studies. 
Here, $V$ is the set of vertices corresponding to a table segment or a passage, and $E$ is the set of edges representing relationships between (table segment, passage) pairs. 
The mapping $\Phi: V \rightarrow \{\texttt{table}\ \texttt{segment}, \texttt{passage}\}$ maps each node to its type, while $\Gamma$ maps a node to its attributes, such as the text of a passage or the matrix of table structures. 
The mapping $\Lambda: E \rightarrow \mathbb{R}$ maps each edge to its score.

The corpus can be expressed as the initial graph $G_{init} = (V_{init}, \emptyset, \Phi, \Gamma, \Lambda_{init})$, where each node in $V_{init}$ one-to-one corresponds to a table segment or passage in $\mathcal{C}$.  
Early fusion generates table-text relationships via entity linking and updates $G_{init}$ before a query $q$ is presented.  
Given $q$, late fusion dynamically generates query-dependent table-text relationships to update $G_{init}$.  
Finally, we retrieve a query-relevant edge-scored bipartite graph $G_q = (V_q, E_q, \Phi, \Gamma, \Lambda_q)$ from $G_{init}$.  
This problem is often interpreted as ranking edges $\mathcal{E}_q$ from all possible edges, as retrieved results are fed to a reader with limited context size~\citep{CORE, COS}.  
$\mathcal{E}_q$ is typically obtained by sorting each edge $e$ in $G$ using its edge scores $\Lambda_q(e)$.  

% The corpus can be expressed as the initial graph $G_{init} = (V_{init}, \emptyset, \Phi, \Gamma, \Lambda_{init})$, where each node in $V_{init}$ one-to-one corresponds to a table segment or a passage in $\mathcal{C}$. 
% Early fusion generates table-text relationships via entity linking and updates $G_{init}$ before a query $q$ is presented. 
% Given a query $q$, late fusion dynamically generates query-dependent table-text relationships to update $G_{init}$. 
% Finally, we aim to retrieve a query-relevant edge-scored bipartite graph $G_q = (V_q, E_q, \Phi, \Gamma, \Lambda_q)$ from $G_{init}$. 
% This problem is often interpreted as finding a ranked sequence of edges $\mathcal{E}_q$ from all possible edges, as the retrieved results are fed to a reader with limited context size~\citep{CORE, COS}.
% $\mathcal{E}_q$ is often generated by sorting each edge $e$ in $G$ using its edge scores $\Lambda(e)$.